% chapters/16/section-04.tex
\section{Lembar dan Pelaksanaan \istilah{Playtest}}
\label{sec:lembar-playtest}

\subsection{Mengamati, Bukan Bertanya}

\istilah{Playtest} yang dilakukan dengan benar menghasilkan keterangan yang
berbeda dari yang Anda duga. Pertanyaan ``apakah permainannya bagus'' tidak
berguna, sebab jawabannya sopan dan tidak dapat ditindaklanjuti. Yang berguna
adalah pengamatan atas apa yang benar-benar dilakukan penguji.

Perbedaan keduanya perlu Anda pahami sebelum menyusun lembar
\istilah{playtest}. Ketika penguji mengatakan permainannya ``cukup jelas'',
Anda tidak memperoleh apa pun. Ketika Anda melihat penguji berhenti selama
dua puluh detik di depan pintu yang belum terbuka, lalu berjalan ke arah yang
salah, Anda memperoleh satu temuan yang dapat langsung Anda perbaiki.

\begin{tipsbox}
Catat waktu, bukan hanya kejadian. ``Penguji berhenti 20 detik di depan pintu
sebelum menemukan jalannya'' jauh lebih berguna daripada ``penguji sempat
bingung''. Angka waktu memudahkan Anda membandingkan hasil perbaikan Anda
dengan hasil sebelumnya.
\end{tipsbox}

\subsection{Isi Lembar \istilah{Playtest}}

Lembar \istilah{playtest} memuat dua bagian: keterangan tentang penguji, dan
catatan atas apa yang diamati selama permainan berlangsung. Bentuknya
diperlihatkan pada Tabel~\ref{tab:lembar-playtest}.

\begin{table}[htbp]
\centering
\caption{Lembar \istilah{playtest} yang dipakai pada praktikum bab ini}
\label{tab:lembar-playtest}
{\setlength{\tabcolsep}{4pt}\footnotesize%
\begin{tabularx}{\textwidth}{@{}L{4.2cm} Y@{}}
\toprule
\headrow \textbf{Bagian} & \textbf{Isi} \\
\midrule
Keterangan penguji & Nama atau kelompok, dan pengalaman bermain permainan
sejenis. \\
\addlinespace
Berkas yang diuji & Nama \istilah{build} beserta tanggalnya. \\
\addlinespace
Waktu penyelesaian & Lama permainan dari awal sampai tujuan, atau keterangan
bahwa pemain berhenti sebelum selesai. \\
\addlinespace
Titik berhenti & Tempat dan waktu ketika penguji berhenti lebih lama dari
sepuluh detik, beserta apa yang tampak ia cari. \\
\addlinespace
Kesalahan yang terlihat & Setiap keadaan yang tidak seharusnya terjadi,
beserta urutan tindakan yang menyebabkannya. \\
\addlinespace
Masukan penguji & Pendapat penguji, dicatat apa adanya, tanpa dibantah. \\
\addlinespace
Catatan penguji sendiri & Hal yang Anda amati tetapi tidak disebutkan
penguji. \\
\bottomrule
\end{tabularx}}
\end{table}

\subsection{Pelaksanaan \istilah{Playtest} antar Kelompok}

\istilah{Playtest} pada praktikum bab ini dilakukan antar kelompok, dan
keuntungannya berjalan dua arah: kelompok Anda memperoleh penguji yang belum
pernah melihat karya Anda, dan Anda memperoleh pengalaman menguji karya orang
lain.

\begin{langkah}
  \item \textbf{Siapkan \istilah{build} yang dapat dijalankan} pada komputer
        lain, bukan pada komputer Anda. Berkas yang hanya berjalan di komputer
        pembuatnya bukan \istilah{build}.
  \item \textbf{Tuliskan tujuan sesi} pada lembar \istilah{playtest}: bagian
        mana yang paling Anda ragukan. Pengujian yang diarahkan akan
        menemukan lebih banyak daripada pengujian yang dibiarkan bebas.
  \item \textbf{Serahkan kendali tanpa memberi petunjuk.} Perkenalkan hanya
        tujuan permainannya, lalu diam. Jangan menjawab pertanyaan dengan
        jawaban; jawablah dengan ``apa yang Anda duga seharusnya terjadi''.
  \item \textbf{Amati dan catat}, termasuk pada menit-menit pertama. Menit
        pertama adalah menit yang paling banyak memuat temuan, dan menit itu
        pula yang paling cepat Anda lupakan.
  \item \textbf{Kumpulkan kesan sesudahnya}, bukan selama permainan.
        Menghentikan permainan di tengah jalan untuk bertanya merusak
        pengalamannya.
  \item \textbf{Tukar peran.} Uji karya kelompok lain dengan cara yang sama,
        agar Anda memahami bagaimana rasanya diamati.
  \item \textbf{Ringkas temuan segera}, selagi ingatan Anda masih segar, lalu
        pindahkan temuan yang dapat diulang ke log \istilah{bug}.
\end{langkah}

\begin{peringatanbox}
Jangan membela karya Anda selama sesi berlangsung. Penguji yang merasa
pendapatnya diperdebatkan akan berhenti memberi masukan, dan sesi berikutnya
menjadi sia-sia. Bantahan, apabila memang perlu, disampaikan sesudah sesinya
selesai, dan biasanya tidak perlu sama sekali.
\end{peringatanbox}
